\section*{Complete author proof (author source pp.34--35)}
The proof depends on the following lemma.
\begin{originalresult}{Lemma 5.18}
For any function $f:\bits^n\to\bits$,
\begin{align*}
\DT(f)&\overset{(1)}\le2\deg(f)^4\overset{(2)}\le O((\log\operatorname{spar}(f\comp\OR))^4)\\
&\overset{(3)}=O((\log\operatorname{rank}((f\comp\OR)_\oplus))^4)\overset{(4)}\le O(D((f\comp\OR)_\oplus)^4)\\
&\overset{(5)}\le O(\DT_\oplus(f\comp\OR)^4).
\end{align*}
\end{originalresult}
\begin{proof}
\begin{enumerate}
\item[(1)] See Buhrman \& de Wolf (2002).
\item[(2)] We need to prove that
\[\deg(f)\le\log\operatorname{spar}(f\comp\OR).\]
We can assume that $f:\{-1,1\}^n\to\{-1,1\}$ has degree $n$. If not, we can fix a monomial of the maximum degree and substitute 1 into all other variables. This substitution does not change the degree, while the sparsity can only decrease. So, the function $f$ has the following multilinear representation:
\[f(x)=\alpha_{[n]}\prod_{i=1}^n x_i+\sum_{S\subsetneq[n]}\alpha_Sx_S,\qquad\text{where }\alpha_{[n]}\ne0.\]
In order to get the multilinear representation of the function $f\comp\OR$ one can substitute each $x_i$ with $\frac{y_{i1}y_{i2}-y_{i1}-y_{i2}-1}{2}$. After the substitution, for any combination of indices $j_1,j_2,\ldots,j_n\in\{1,2\}$, the coefficient of the monomial $\prod_{k=1}^n y_{kj_k}$ is equal to $\alpha_{[n]}\cdot2^{-n}$. Indeed, such monomials can only be produced by the substitution into the maximum degree monomial, and such substitution results in the coefficient $\alpha_{[n]}\cdot2^{-n}$.
Note that $\alpha_{[n]}\ne0$ and hence there is a monomial for every $2^n$ possible combinations of indexes $j_1,j_2,\ldots,j_n$. That gives us the desired bound:
\[\log\operatorname{spar}(f\comp\OR)\ge n=\deg(f).\]
\item[(3)] See Bernasconi \& Codenotti (1999).
\item[(4)] See Lov\'asz \& Saks (1993).
\item[(5)] Each time we query a linear function $\ell(x\oplus y)=\ell(x)\oplus\ell(y)$, Alice and Bob can send to each other $\ell(x)$ and $\ell(y)$ to compute $\ell(x\oplus y)$.
\end{enumerate}
\end{proof}
Theorem 5.16 follows immediately. For Theorem 5.17, we need one more inequality in the chain.
\begin{proof}[Theorem 5.17]
We need to show that
\[\DT_\oplus(f\comp\OR)\le2\DT(f).\]
Indeed, the decision tree for $f$ can be transformed into a decision tree for $f\comp\OR$ by replacing every query with two queries of corresponding variables. Together with Lemma 5.18 that finishes the proof.
\end{proof}
